% !TeX root = test-main-adv.tex

\begin{center}
{\Large\textbf{最終テスト}}
\end{center}
\vspace{0.15em}
\begin{flushright}
氏名：\testnameblank
\end{flushright}
\vspace{0.25em}
% ===== リスニング（35点）解答欄のみ =====
{\underline{\large\noindent\textbf{リスニング} ：問題用紙・音声は別紙。下記の解答欄に記入すること。（35点）}}
\vspace{0.5em}

\needspace{5\baselineskip}
\noindent\textbf{問1}
\vspace{0.4em}

\noindent
No.\,1 \listeningchoice[ア]{3.2em}（2点）\qquad
No.\,2 \listeningchoice[ウ]{3.2em}（2点）\qquad
No.\,3 \listeningchoice[ア]{3.2em}（2点）
\testqseplistening

\needspace{6\baselineskip}
\noindent\textbf{問2}
\vspace{0.4em}

\noindent
No.\,1 \listeningchoice[エ]{3.2em}（3点）\qquad
No.\,2 \listeningchoice[イ]{3.2em}（3点）

\vspace{0.75em}

\noindent
No.\,3 \listeningchoice[ア]{3.2em}（3点）\qquad
No.\,4 \listeningchoice[ウ]{3.2em}（3点）
\testqseplistening

\needspace{5\baselineskip}
\noindent\textbf{問3}
\vspace{0.4em}

\noindent
No.\,1 \listeningchoice[イ]{3.2em}（3点）\qquad
No.\,2 \listeningchoice[ウ]{3.2em}（3点）\qquad
No.\,3 \listeningchoice[エ]{3.2em}（3点）
\testqseplistening

\needspace{7\baselineskip}
\noindent\textbf{問4}
\vspace{0.4em}

\noindent No.\,1\quad He enjoyed it in a tent on \listeningblank[the lake]{0.5\linewidth} （4点）

\vspace{0.85em}

\noindent No.\,2\quad He \listeningblank[cooked fish]{0.5\linewidth} and ate them with Tomoya. （4点）

\testsectsep\vspace{1.5em}
% ===== 筆記（65点） =====
\underline{\large\noindent\textbf{筆記} ：以下の問題は、配布テキストに沿って解答すること。（65点）}
\vspace{0.35em}
\testsectsep
% 2点問題シャッフル順（seed 20260804、問15除く）
\testquestion
\noindent\textbf{【問1】} 次の日本語を主語と動詞を含む英文にしなさい。

\noindent「彼は疲れているにちがいない。」

\testblank[He must be tired.]{0.92\linewidth}
\testqsep
\testquestion
\noindent\textbf{【問2】} 次の英文を意味を変えないように書き換えなさい。

\noindent「No other city in Japan is as large as Tokyo.」

\testblank[Tokyo is the largest city in Japan.]{0.92\linewidth}
\testqsep
\testquestion
\noindent\textbf{【問3】} 次の英文を疑問文にしなさい。

\noindent「He can play the guitar very well.」

\testblank[Can he play the guitar very well?]{0.92\linewidth}
\testqsep
\testquestion
\noindent\textbf{【問4】} 次の英文に対する答えを主語と動詞を含む英文1文で答えなさい。

\noindent「Which country has the largest population in the world?」

\testblank[India has the largest population in the world.]{0.92\linewidth}
\testqsep
\testquestion
\noindent\textbf{【問5】} 次の英文に対する答えを主語と動詞を含む英文で答えなさい。

\noindent「Would you like a cup of coffee?」

\testblank[Yes, please. / No, thank you.]{0.92\linewidth}
\testqsep
\testquestion
\noindent\textbf{【問6】} 次の英文に主語と動詞を含む1文で答えなさい。

\noindent「Why are you happy?」

\testblank[Because I passed the exam.]{0.92\linewidth}
\testqsep
\testquestion
\noindent\textbf{【問7】} 次の英文を和訳しなさい。

\noindent「This room is not as large as that one.」

\testblank[この部屋はあの部屋ほど広くない。]{0.92\linewidth}
\testqsep
\testquestion
\noindent\textbf{【問8】} 次の英文を和訳しなさい。

\noindent「My bag is heavier than yours.」

\testblank[私のかばんはあなたのかばんより重い。]{0.92\linewidth}
\testqsep
\testquestion
\noindent\textbf{【問9】} whoが主語になるときの疑問文を作りなさい。また、\textbf{その疑問文に答えなさい}。

\vspace{0.5em}
\testblank[Who is the teacher? \quad Mr. Riku is.]{0.92\linewidth}
\testqsep
\testquestion
\noindent\textbf{【問10】} 次の英文を否定文にしなさい。

\noindent「She will visit Kyoto next week.」

\testblank[She won't visit Kyoto next week.]{0.92\linewidth}
\testqsep
\testquestion
\noindent\textbf{【問11】} 次の英文を意味を変えないように書き換えなさい。

\noindent「Mike is taller than Tom.」

\testblank[Tom is not as tall as Mike.]{0.92\linewidth}
\testqsep
\testquestion
\noindent\textbf{【問12】} 次の日本語を、主語と動詞を含む英文1文で表現しなさい。

\noindent「彼女は、私を幸せにしてくれます。」

\testblank[She makes me happy.]{0.92\linewidth}
\testqsep
\testquestion
\noindent\textbf{【問13】} whatが主語になるときの疑問文を作りなさい。また、\textbf{その疑問文に答えなさい}。

\vspace{0.5em}
\testblank[What happened here? \quad A big accident happened.]{0.92\linewidth}
\testqsep
\testquestion
\noindent\textbf{【問14】} 次の日本語を英語に直しなさい。

\noindent「メアリーに代わっていただけますか?」

\testblank[May I speak to Mary?]{0.92\linewidth}
\testqsep
\testquestion
\noindent\textbf{【問15】} 次の日本語を、主語と動詞を含む英文1文で表現しなさい。

\noindent「コミュニケーション能力を高めたい。」

\testblank[I want to talk with many people well.]{0.92\linewidth}
\testqsep

\newpage
\testquestion
\noindent\textbf{【問16】} 次の慣用表現を英語に直しなさい。

\noindent「(道案内で)札幌駅までの行き方を教えていただけますか?」

\testblank[Could you tell me how to get to Sapporo Station?]{0.92\linewidth}
\testqsep
\testquestion
\noindent\textbf{【問17】} 次の日本語を主語と動詞を含む英文にしなさい。

\noindent「あなたのペンを借りてもいいですか。」

\testblank[May I use your pen?]{0.92\linewidth}
\testqsep
\testquestion
\noindent\textbf{【問18】} 次の英語を日本語に直しなさい。

\noindent「What would you like to order?」

\testblank[何をご注文でしょうか?]{0.92\linewidth}
\testqsep
\testquestion
\noindent\textbf{【問19】} 次の英文を和訳しなさい。

\noindent「What a cool math teacher he is!」

\testblank[なんてカッコいい数学の先生でしょう。]{0.92\linewidth}
\testqsep

\testsectsep
\testquestionwide
\noindent\textbf{【問20】} 次の英文に対する答えを主語と動詞を含む英文1文で自由に書きなさい。（5点）

\noindent「What do you think about reading English books? And why?」

\testblank[I think it is good because it helps me learn new words.]{0.92\linewidth}
\testqsep
\testquestionwide
\noindent\textbf{【問21】} 自分が将来、できるであろうことについて、主語と動詞を含む英文1文を作りなさい。（5点）

\testblank[I will be able to speak English well.]{0.92\linewidth}
\testqsep
\testquestionwide
\noindent\textbf{【問22】} 次の英文に対する答えを主語と動詞を含む英文1文で自由に書きなさい。（5点）

\noindent（選択肢：Sapporo, New York, London）「Which of the three cities do you want to go to the most? And why?」

\vspace{0.5em}
\testblankwide[I want to go to New York because it is the biggest city in the world.]{0.92\linewidth}
\testqsep

\newpage
\testquestionwide
\noindent\textbf{【問23】} 次の英文は、ある中学生が、英語の授業で、「私の人生で最も大きな経験」について書いたものです。あなたがその中学生になったつもりで、条件にしたがって、英文を完成させなさい。（12点）

\vspace{0.3em}
\noindent\textbf{英文}
\vspace{0.15em}

\noindent\fbox{%
\parbox[t]{0.97\linewidth}{%
\vspace{0.35em}%
\noindent
The \compnumblank{1} experience in my life is \compnumblank{2}.\\[0.45em]
\noindent
\underline{There are two reasons for this.}\ \compnumblank{3}%
\vspace{0.35em}%
}}%

\vspace{0.5em}
\noindent\textbf{条件}
\vspace{0.2em}

\noindent\doublebox{%
\parbox[t]{0.95\linewidth}{%
\vspace{0.25em}%
\noindent ・\,(1), (2) には、英文に合わせて適当な英語を書きなさい。\\[0.25em]
\noindent ・\,(3) には、下線部\underline{\hspace{3em}}の内容につながるよう、24語以上の英語で書きなさい。%
\vspace{0.25em}%
}}%

\vspace{0.55em}
\noindent The \compnumblank[biggest]{1} experience in my life

\vspace{0.45em}
\noindent is \compnumblankwide[studying English hard every day]{2} .

\vspace{0.55em}
\noindent\textbf{(3)}

\vspace{0.35em}
\wordentryrowsix{6}{First,}{I}{made}{many}{new}{friends.}
\wordentryrowsix{12}{in}{the}{club.}{Second,}{I}{became}
\wordentryrowsix{18}{more}{interested}{in}{foreign}{countries.}{I}
\wordentryrowsix{24}{want}{to}{visit}{the}{U.S.A.}{someday.}
\wordentryrowsix{30}{}{}{}{}{}{}
\wordentryrowsix{36}{}{}{}{}{}{}

\vfill
\begin{flushright}
リスニング \underline{\hspace{8em}} ／ 35点\qquad
筆記 \underline{\hspace{8em}} ／ 65点\\[1.5em]
\textbf{合計 \underline{\hspace{8em}} ／ 100点}
\end{flushright}
